\documentclass[conference]{IEEEtran}
\IEEEoverridecommandlockouts
\usepackage{cite}
\usepackage{amsmath,amssymb,amsfonts}
\usepackage{graphicx}
\usepackage{booktabs}
\usepackage{textcomp}
\usepackage{xcolor}
\usepackage{url}
\usepackage{pgfplots}
\pgfplotsset{compat=1.18}
\def\BibTeX{{\rm B\kern-.05em{\sc i\kern-.025em b}\kern-.08em
    T\kern-.1667em\lower.7ex\hbox{\sc e}\kern-.125em X}}

\begin{document}

\title{Signal Diagnostics and a Fourteen-Lane Evaluation Ledger\\
for GRPO Post-Training: A Provenance-Bound Measurement Report}

\author{\IEEEauthorblockN{Arvind C R}
\IEEEauthorblockA{\textit{Dept.\ of Computer Science and Engineering} \\
\textit{PES University}\\
Bengaluru, India \\
arvindcr4@gmail.com}
\and
\IEEEauthorblockN{Ramesh Prakash Guledgudd}
\IEEEauthorblockA{\textit{Dept.\ of Computer Science and Engineering} \\
\textit{PES University}\\
Bengaluru, India}
}

\maketitle

\begin{abstract}
This is a measurement and audit report, not a new estimator or a capability
claim. We report short-horizon GRPO-family telemetry and a governance case
study: a fourteen-lane evaluation ledger (E1--E14) for one short-horizon
actor (80 prompt draws over 40 steps; executed training rows unrecovered)
whose paired rerun found no difference from its base model, i.e.\ ledger
mechanics demonstrated on a null, partly undisclosed intervention. Uniform
rewards remove the centred reward-driven advantage term; the Zero-Variance
Fraction (ZVF) counts such groups and is bookkeeping over logged rewards,
not a validated predictor of learning, gradient magnitude or capability.
Stored tensors and logged runs come from different, uncrossed populations
and are reported separately; no effect is allocated to stack versus
algorithm. Same-stack algorithm contrasts are inconclusive: PPO versus GRPO
over five paired seeds gives $\Delta=-0.002$ (paired $t_4=-1.0$,
$p$-value $=0.374$, at a 0.98--1.00 ceiling), and four GRPO variants over
eight paired seeds give no contrast detectable at their minimum detectable
effects. The campaign records 129/312 (41.35\%) on VerilogEval and
2,271/4,426 (51.31\%) among judged Omni-MATH rows, but only three lane
scopes are strictly complete and two further scopes carry terminal or
denominator qualifications (Section~IV-F taxonomy); no full-suite
base-model comparator was run. No training benefit, cross-suite aggregate,
algorithm equivalence or independent stack attribution is established.
\end{abstract}

\begin{IEEEkeywords}
reinforcement learning, GRPO, large language models, verifiable rewards,
reproducibility, benchmark evaluation
\end{IEEEkeywords}

\section{Introduction}

Reinforcement learning with verifiable rewards (RLVR) replaces the learned
reward model of RLHF~\cite{ouyang2022training} with a programmatic checker,
such as exact match for mathematics or unit tests for code. The workhorse
optimiser was long PPO~\cite{schulman2017proximal}, which needs a value
network. GRPO~\cite{shao2024deepseekmath,guo2025deepseekr1} removes the critic:
it samples a group of $G$ completions per prompt and standardises each reward
against the group's own statistics. For one group $g$ with rewards
$r_1,\dots,r_G$ and a small stability constant $\varepsilon>0$,
\begin{equation}
A_i = \frac{r_i - \mu_g}{\sigma_g + \varepsilon}, \;
\mu_g = \tfrac{1}{G}\textstyle\sum_{j=1}^{G} r_j, \;
\sigma_g^2 = \tfrac{1}{G}\textstyle\sum_{j=1}^{G} (r_j-\mu_g)^2 ,
\label{eq:adv}
\end{equation}
for $i\in\{1,\dots,G\}$. Implementations differ in this normalisation
($1/G$ versus $1/(G-1)$, group versus batch standard deviation, or none).

This simplification has two consequences that standard training logs do not
reveal. The first concerns \emph{mechanism}. When every completion in a group
receives the same reward, the numerator $r_i-\mu_g$ is zero for every $i$
(while $\varepsilon$ keeps the denominator positive), so every $A_i$ in
(\ref{eq:adv}) is exactly zero; the same holds without $\sigma_g$ scaling.
The group has been sampled and scored at full cost, but its centred
reward-driven advantage term is zero. This statement is scoped to that term:
deployed variants may still apply KL, entropy, clipping or length terms, and
batch-level normalisation can couple groups, so ``uniform group'' does not
mean ``zero total gradient'' under every estimator (see the per-estimator
scoping in Section~III-A). The second concerns
\emph{provenance}. ``GRPO'' is implemented in Tinker~\cite{thinkingmachines2024tinker},
TRL~\cite{vonwerra2022trl}, veRL~\cite{sheng2024verl},
OpenRLHF~\cite{hu2024openrlhf} and many forks, each with its own answer
parser, temperature, KL coefficient, clip range, LoRA targets~\cite{hu2022lora}
and evaluator. Two runs labelled GRPO can differ in many ways that plausibly
change the outcome.

This paper is a measurement/audit report from an individual M.Tech project.
It proposes no new estimator and establishes no training effect. Our
contributions, scoped to the delivered evidence, are:
\begin{itemize}
\item an audit methodology: source-bound recomputation of ZVF bookkeeping
from stored reward tensors where they survive, with logged-only material
kept descriptive and never pooled with tensor-backed material; the dispatch
harness is infrastructure for this audit (completed runs come only from
Tinker and TRL; veRL and OpenRLHF entries are unexecuted launch plans);
\item ZVF as a logging/bookkeeping standard with its elementary tie-probability
baseline, plus an offline arithmetic replay of an adaptive group-size
trigger (no controller benefit is claimed; predictive validity of ZVF for
held-out learning outcomes is unestablished);
\item inconclusive same-stack contrasts and label-collision exhibits,
reported with seed-level statistics and minimum detectable effects, and an
eight-item minimum reportable stack proposed as a reporting standard;
\item a fourteen-lane evaluation ledger as a governance case study on one
null, partly undisclosed actor (item-level training separation not
established, so the campaign is not claimed as a verified held-out
experiment), with each lane reported at its terminal state under a single
completeness taxonomy.
\end{itemize}

\section{Related Work}

GRPO was introduced with DeepSeekMath~\cite{shao2024deepseekmath} and
popularised by DeepSeek-R1~\cite{guo2025deepseekr1}. Variants modify the
estimator: DAPO~\cite{yu2025dapo} adds decoupled clipping and dynamic sampling
that discards zero-variance groups, GSPO~\cite{qwen2025gspo} moves the ratio to
the sequence level, and Dr.\ GRPO~\cite{liu2025understanding} removes the
length and standard-deviation normalisations it identifies as biased.
RLOO~\cite{ahmadian2024backtobasics} uses a leave-one-out baseline, which for
binary groups is proportional to the centred advantage and vanishes on the
same uniform groups. Zero-variance groups are already targeted by prior work:
GRESO~\cite{zheng2025greso} predicts and skips them before rollout,
AERO~\cite{zhang2026aero} adapts rollout allocation in group-based RL, and
RL-ZVP~\cite{le2025rlzvp} extracts signal from zero-variance prompts by
entropy-guided advantage shaping. The tie probability
$p^G+(1-p)^G$ for i.i.d.\ binary rewards is elementary; we use it as a
measurement baseline, not as a new result. Preference-based alternatives such
as DPO~\cite{rafailov2024direct} and open recipes such as
T\"ULU~3~\cite{lambert2024tulu3} complete the landscape. On methodology,
Henderson et al.~\cite{henderson2018deep} and Engstrom et
al.~\cite{engstrom2020implementation} showed that deep RL results are
sensitive to seeds and code-level details; reproducibility
checklists~\cite{pineau2020improving}, model cards~\cite{mitchell2019modelcards},
datasheets~\cite{gebru2021datasheets} and evaluation
frameworks~\cite{liang2022helm,biderman2024lmeval} set reporting norms. We do
not propose a new estimator; the uncrossed corpus cannot allocate effects to
stack versus algorithm.

\section{Methodology}

\subsection{Provenance-bound design}
Matched comparisons aim to hold model, prompts, verifier, decoding configuration,
step budget and evaluation fixed while changing one factor; the historical
cross-backend summaries below do not meet that design, and no task in this
corpus is replicated across two independent stacks with tensors intact, so
manifest-bound entries may differ without identifying a stack effect
(confounding is tabulated in Section~IV-D). A shared task,
reward and decoding specification is dispatched through thin adapters to each
back-end, so the same experiment runs without re-implementing the grader.
Being able to run the same experiment on every back-end does not mean every
back-end completed it. In this release only Tinker and TRL produced completed
runs; veRL and OpenRLHF entries are dry-run launch plans and are never pooled
with real results. Per-estimator scoping of the uniform-group statement:
under (\ref{eq:adv}) and under mean-centering without $\sigma_g$ division
(the runner that produced the Phase-1 tensors), a uniform group zeroes the
centred advantage for its members; DAPO-style dynamic sampling instead
discards such groups, Dr.\ GRPO removes the $\sigma$ scaling (the centred
term still vanishes), and sequence-level GSPO changes the ratio level; in
all cases KL/entropy/length terms, if enabled, are outside the ZVF
bookkeeping reported here. Every ZVF number below is scoped to the estimator
and runner that produced it.

\subsection{Telemetry and the Zero-Variance Fraction}
The training loop emits per-step ZVF, gradient utilisation $\mathrm{GU}=1-\mathrm{ZVF}$,
policy entropy, completion length and KL divergence; the tensor-instrumented
runners also persist the raw per-group reward tensors. For a batch
$\mathcal{B}$ of $|\mathcal{B}|$ prompt groups,
\begin{equation}
\mathrm{ZVF} = \frac{1}{|\mathcal{B}|}\sum_{g\in\mathcal{B}}
\mathbf{1}\left[\sigma_g = 0\right].
\label{eq:zvf}
\end{equation}
Where tensors are stored, ZVF is recomputed exactly from them.
$\mathrm{GU}=1-\mathrm{ZVF}$ is a rename of the complement, used only as a
plotting axis, not a separate metric. ZVF is unsigned and requires mean
reward for interpretation; it does not distinguish saturated mastery,
impossible prompts, reward coarseness, truncation or checker failure.

\subsection{A tie-probability baseline for group size}
Assume a prompt with success probability $p$ yields $G$ conditionally i.i.d.\
binary rewards $r_j\sim\mathrm{Bernoulli}(p)$. Then the expected ZVF of its
group is the tie probability $h_G(p) = p^G + (1-p)^G$, which alone implies
(i) a U-shape of ZVF in $p$, with minimum $2^{1-G}$ at $p=0.5$, and (ii) a
strict decrease in ZVF with $G$ for $p\in(0,1)$. These are restatements of
the algebra, not corroborating tests: agreement of binned ZVF with
$h_G(\hat p)$ computed from the same rewards is bookkeeping consistency,
not out-of-sample prediction. Mixtures of prompt
difficulties within a batch raise batch ZVF above $h_G(\bar p)$; deviations
from $h_G(\bar p)$ therefore encode prompt-mixture dispersion and possible
within-group correlation (shared prefix, sampling temperature, checker
determinism), which we do not separately estimate here, so ZVF cannot
distinguish ``prompt too easy/hard'' from ``sampling too correlated.'' As a
plotting heuristic only, with no predictive claim, we show
\begin{equation}
S(p, G) = p(1-p)\,\bigl(1 - h_G(p)\bigr),
\label{eq:signal}
\end{equation}
an inverted-U in $p$. The $p(1-p)$ factor is the variance of an
\emph{unnormalised} centred advantage; under the $\sigma_g$ scaling of
(\ref{eq:adv}) the mean squared advantage of a non-uniform group is
approximately one, so (\ref{eq:signal}) is not an expected normalised-gradient
identity and both of its factors peak at $p=0.5$ by construction. No
gradient-norm prediction is tested in this paper (see Section~IV-B);
predictive validity of ZVF or $S$ for future learning outcomes is
unestablished. The adaptive trigger replayed arithmetically in Section~IV-C
raises $G$ for a prompt when its observed group is uniform.

\subsection{Statistical protocol}
The seed is the unit of replication. Matched contrasts use two-sided paired
$t$-tests on per-seed differences at $\alpha=0.05$, with degrees of freedom
reported; minimum detectable effects (MDE) are for 80\% power from the
noncentral $t$; the four variant-versus-GRPO contrasts of Section~IV-D form
one Benjamini--Hochberg family. Consecutive training steps are not independent
observations. Single-seed panels, logged run summaries and offline trace audits
are reported descriptively without inferential claims. Inconclusive tests do
not establish equivalence.

\subsection{Evaluation campaign (E1--E14): a governance case study}
The campaign is reported as ledger mechanics on a null, partly undisclosed
actor, not as a benchmark result for a validated training intervention. It
evaluates one actor, Qwen3.6-35B-A3B with a Tinker-trained LoRA
adapter (rank 32, all linear layers) in bfloat16. Its recovered training
record reports 40 steps at $G=4$ with 2 prompts per step (80 prompt draws)
from a 512-row API-Bank/SWE-Gym tool-use pool, learning rate $10^{-5}$,
temperature 0.7 and seed 809; reward was zero on 22 of 40 steps. The exact
executed training code, the 80 drawn prompts and the consumed-row inventory
are not recovered, so item-level training separation is not established:
the campaign is therefore not described as held-out in the blanket sense,
and it is not a verified training-effect experiment. The 512-row pool
composition and any decontamination check against the evaluation lanes
cannot be published from recovered records (AUTHOR\_TODO: supply the pool
manifest and overlap check, or retain this non-claim). The primary
training-effect statement is the small paired trained-versus-base rerun
(Section~V), which found no difference. Each lane names one
original benchmark contract; several also have a declared replacement scope
(Table~\ref{tab:lanes}). Original and replacement scopes are never pooled,
and replacement executions do not close their original contracts. Campaign
timeline: original contracts and terminal states were frozen in the
19~September ledger; six replacement executions under a protocol override
followed on 27~September (marked exploratory), and the paired rerun was run
on 26~September; lane-provenance sources were accessed 18~October 2026.

\section{Results}

\subsection{Zero-variance frequency in the Phase-1 configuration}
The Phase-1 configuration is Qwen3.5-4B on GSM8K~\cite{cobbe2021training} on
Tinker, $G=8$, 16 prompts per step, 40 steps, one seed, four runner labels
(GRPO, AERO, GIFT, AReaL) in a mean-centred runner without $\sigma_g$
division, so these measurements are scoped to that estimator, not to
``GRPO'' generally. The labels are internal runner dispatches, defined here
at first use from the dispatch code
(\texttt{variance\_mitigation\_integration.py}): GRPO is the unmodified
mean-centred baseline; AERO sets an adaptive-group-size flag (grow/shrink
$G$ on rolling ZVF thresholds; cf.\ the external method of the same
name~\cite{zhang2026aero}, of which this flag is not a verified
reimplementation); GIFT sets an implicit-reward-matching flag in the
advantage computation; AReaL sets an asynchronous-training/stale-buffer
flag. GIFT and AReaL have no archival citation in this paper and must not be
read as the external systems of those names (AUTHOR\_TODO: supply archival
definitions or rename the internal flags before submission).
Hash comparison of the stored tensors (Appendix~A) shows the four files are
byte-distinct and GRPO/AERO rewards differ in 225 of 640 groups, yet every
reported GRPO and AERO summary coincides exactly (mean step ZVF 0.7203;
461 degenerate groups; 427 all-correct; 34 all-wrong), and $G$ remained 8 in
all 40 AERO steps, so the adaptive dispatch never manifested. The panel is
therefore reported as four labels with at most three distinct summary
profiles, the GRPO/AERO coincidence is disclosed rather than read as a
cross-method comparison, and no cross-method conclusion is drawn from
Section~IV-A or IV-C. Recomputed from the stored tensors (640 groups per label), mean
step ZVF is 0.706--0.770 and all-correct groups are 0.647--0.711 of all
groups, so most uniform groups are prompts already solved in the observed
group; this does not establish true $p=1$. The regime is configuration-specific:
a separate sampling-only cohort (Qwen3-8B on 200 GSM8K test problems per seed,
$G=8$, temperature 1.0, three seeds, no training) has interior accuracy
(0.67--0.71) and stored-tensor ZVF of 0.130, 0.190 and 0.155 (pooled 0.158).
It is a different population and is not compared with the Phase-1 panel.

\subsection{Group size: a mechanism, not a universal optimum}
In a 12-run sweep (Qwen2.5-0.5B, synthetic arithmetic task, $G\in\{2,4,8,16\}$,
three seeds, 40 steps) mean ZVF falls from 0.838 at $G=2$ to 0.631 at $G=16$
(Fig.~\ref{fig:gu}). Accuracy on the split reserved from training is at ceiling (0.97--0.995 per seed),
so the sweep establishes no accuracy optimum, and larger $G$ costs
proportionally more rollouts. The logged \texttt{zvf-audit} project (368 W\&B
runs; nine model labels; $G\in\{4,6,8,16\}$; 2--12 steps; 350 finished)
is a heterogeneous offline summary, not a common-protocol cohort. Binned by
final reward over all 350 finished runs (18 crashed/failed runs excluded;
bin edges are post-hoc deciles of final reward, not pre-registered), mean
run ZVF is 0.959 for the 207 zero-reward runs and 0.881 for the 52 runs at
reward $\ge0.9$; the interior evidence rests on only 10 runs with final
reward in $(0.1,0.6)$ (per-decile $n=2,3,1,1,3$; mean ZVF 0.12--0.38), and
the remaining 81 finished runs fall at reward $[0.6,0.9)$ (decile
$n=11,26,44$; mean ZVF 0.822, 0.773, 0.812). The three bins quoted in
earlier drafts thus account for 269 of 368 logged runs. The pattern is the
near-tautological $h_G$ bookkeeping shape (final reward near 0 or 1
mechanically implies mostly uniform groups), not an independent test of
prediction (i). Prediction (ii) appears in some model rows
(Qwen3-32B: 0.33, 0.18, 0.08 at $G=4/8/16$ from 2, 2 and 1 runs) but not in
others (Qwen3-4B-Instruct: 0.89, 0.84, 0.92), and model, task, $G$ and
horizon vary together, so no stack reading is drawn. A gradient-norm
association was once logged as a single scalar ($r=+0.71$) for one
Qwen2.5-0.5B run aggregate with the number of observations unlogged and
steps autocorrelated; it is withdrawn as a numbered prediction and retained
only as an unlogged directional observation that must not be cited as a
correlation estimate. Predictive validity of ZVF for gradient magnitude or
held-out gain is unestablished.

\begin{figure}[t]
\centering
\begin{tikzpicture}
\begin{axis}[width=0.95\columnwidth,height=4.6cm,xmode=log,log basis x=2,
 xtick={2,4,8,16},xticklabels={2,4,8,16},xmin=1.7,xmax=19,ymin=0,ymax=0.45,
 xlabel={Group size $G$},ylabel={$\mathrm{GU}=1-\mathrm{ZVF}$},
 ymajorgrids,grid style={gray!25},tick label style={font=\footnotesize},
 label style={font=\footnotesize}]
\addplot[only marks,mark=square*,mark size=2pt,black,
 error bars/.cd,y dir=both,y explicit]
 coordinates {(2,0.1620) +- (0,0.0126) (4,0.2365) +- (0,0.0157)
              (8,0.3094) +- (0,0.0118) (16,0.3688) +- (0,0.0083)};
\end{axis}
\end{tikzpicture}
\caption{Gradient utilisation versus group size (Qwen2.5-0.5B arithmetic,
three seeds, 40 steps). Points are seed means; bars are $\pm1$ seed SD.
No fitted law or extrapolation is shown, and different $G$ consume different
rollout budgets. Source: \texttt{groupsize\_zvf\_sweep.json}.}
\label{fig:gu}
\end{figure}

\subsection{An offline arithmetic replay of the adaptive trigger}
An earlier four-arm open-trainer comparison of fixed and adaptive group size
has no retained per-arm records and is not reported. We report instead an
offline arithmetic replay over the stored Phase-1 tensors (4 labels $\times$ 640
prompt-steps $=2{,}560$ observations). A symmetric trigger that escalates on
any uniform group fires on 1,867 of them; 1,723 (92.3\%) are all-correct
groups and 144 all-incorrect. Deduplicating the coincident GRPO/AERO
summary profile (Section~IV-A) gives 1,406 fires over 1,920 groups, of which
1,296 (92.2\%) are all-correct and 110 all-incorrect; the share is
essentially unchanged and largely follows from the reward level. No controller executed in this replay and no fresh rollouts were drawn,
so it cannot show whether larger $G$ would restore contrast or what the live
cost would be; plugging an observed $\hat p\in\{0,1\}$ into $h_G$ forces zero
predicted rescue by construction. This replay is arithmetic only: we make
no trigger recommendation, positive or negative. Any asymmetric,
cost-aware trigger remains an untested hypothesis requiring a live,
compute-matched A/B (including a larger-$G$ baseline at equal rollout
budget) before it can be recommended or rejected.

\subsection{Same-stack contrasts and label-collision exhibits}
Historical marginal one-way ANOVAs over a 42-experiment subset are not
reported: the row-level membership of that subset is not recoverable, so its
effect sizes cannot be recomputed or adjusted for confounding. No crossed
stack$\times$algorithm cell exists in this corpus; Table~\ref{tab:confound}
shows that every bundle below varies model, task, $G$ and horizon together,
so cross-bundle differences are manifest-bound case studies and no causal
stack language (``stack drives,'' ``stack effect'') is used. Evidence here
is confined to matched same-stack contrasts and exhibits.

\begin{table}[t]
\caption{Confounding across evidence bundles: stack, model, task, group
size and horizon vary together; no task is replicated across two stacks.}
\label{tab:confound}
\centering
\footnotesize
\setlength{\tabcolsep}{3pt}
\begin{tabular}{@{}lllll@{}}
\toprule
Bundle & Stack & Model & Task & $G$, horizon \\
\midrule
Phase-1 (IV-A/C) & Tinker & Qwen3.5-4B & GSM8K & 8, 40 st. \\
Sweep (IV-B) & TRL-style & Qwen2.5-0.5B & arithmetic & 2--16, 40 st. \\
W\&B audit (IV-B) & mixed & 9 labels & mixed & 4--16, 2--12 st. \\
PPO/GRPO (IV-D) & TRL-style & Qwen2.5-0.5B & arithmetic & ---, 40 st. \\
Variants (IV-D) & TRL & Qwen3-8B & GSM8K & 4, 30 st. \\
E1--E14 actor & Tinker & Qwen3.6-35B-A3B & tool-use pool & 4, 40 st. \\
\bottomrule
\end{tabular}
\end{table}

\emph{PPO versus GRPO (protocol demonstration at ceiling).} On one TRL-style stack (Qwen2.5-0.5B, synthetic
arithmetic, 40 steps, five seeds) the per-seed difference on the split
reserved from training, GRPO$-$PPO,
is $-0.002$ (95\% CI $[-0.008, +0.004]$, $t_4=-1.0$, $p=0.374$), at a
0.98--1.00 ceiling. The interval is narrow only because there is nowhere to
move; the design could not have detected any plausible algorithm effect, so
this contrast is retained only as a demonstration of the paired protocol: a
failure to detect a difference, not evidence of equivalence, and it carries
no stack conclusion. An earlier Qwen3-8B PPO/GRPO comparison is a provenance conflict
(different back-ends; its $n=10$ were autocorrelated steps of one run per arm)
and is not used.
\emph{Four variants versus GRPO.} On one TRL stack (Qwen3-8B, GSM8K, $G=4$,
30 steps, eight paired seeds, 500 reserved items each; per-seed scores and
exact statistics in Table~\ref{tab:variants}, source
\texttt{e1\_confirmatory\_wandb\_export.tsv}, Appendix~A) the means are
0.633 (GRPO), 0.634 (DAPO), 0.638 (GSPO), 0.631 (Dr.\ GRPO) and 0.632 (AERO).
All four contrasts are inconclusive at MDEs (0.010--0.015) at or above the
0.010 design margin; inconclusive does not mean equivalent.

\begin{table}[t]
\caption{Variant contrasts versus GRPO (8 paired seeds, df$=7$). Per-seed
reserved-split scores by seed (11/23/37/53/71/89/107/131): GRPO
.650/.636/.634/.634/.634/.620/.634/.622; DAPO
.648/.624/.628/.636/.634/.636/.634/.632; GSPO
.640/.638/.642/.632/.638/.646/.640/.628; Dr.\ GRPO
.642/.626/.622/.636/.620/.628/.628/.646; AERO
.630/.628/.636/.632/.632/.630/.630/.640.}
\label{tab:variants}
\centering
\footnotesize
\setlength{\tabcolsep}{3pt}
\begin{tabular}{@{}lrrrrr@{}}
\toprule
Contrast & $\Delta$ & 95\% CI & $t_7$ & $p$ (BH) & MDE \\
\midrule
DAPO & $+0.001$ & $[-0.006,0.008]$ & 0.32 & 0.756 (0.858) & 0.010 \\
GSPO & $+0.005$ & $[-0.004,0.014]$ & 1.38 & 0.210 (0.841) & 0.012 \\
Dr.\ GRPO & $-0.002$ & $[-0.013,0.009]$ & $-0.44$ & 0.673 (0.858) & 0.015 \\
AERO & $-0.001$ & $[-0.010,0.009]$ & $-0.19$ & 0.858 (0.858) & 0.013 \\
\bottomrule
\end{tabular}
\end{table}
\emph{Label collisions.} Two runs labelled ``DAPO'' on different stacks showed
mean ZVF of 0.00 and 0.58; the managed run's closed internals prevent
verifying the mechanism. Nominally matched Tinker and TRL GRPO runs on
Qwen3-8B GSM8K reached last-10 training rewards of 0.856 and 0.050, but the
base checkpoint also changed (Qwen3-8B-Base versus Qwen3-8B), so this is an
under-specification exhibit, not a back-end effect. On one managed platform
four labels landed within a 0.034 training-reward band (0.710--0.744, three
seeds, 15 steps, no held-out metric); this is not evidence of equivalence.
These exhibits motivate an eight-item minimum reportable stack: (1) loss form,
(2) reference policy and KL handling, (3) sampler and back-end, (4) telemetry,
(5) group-size schedule, (6) held-out split, (7) decontamination, and (8)
held-out pass@$k$ under a stated evaluator. A machine-readable registry diff
flags the ``DAPO'' pair from manifests alone.

\subsection{Fourteen-lane evaluation ledger (governance case study)}
Table~\ref{tab:lanes} gives each lane's original contract, replacement
scope and terminal state. One taxonomy is used everywhere in this paper
(abstract, text and tables): \emph{strictly complete} (full scope-specific
dispositions), \emph{terminal-complete} (complete with a stated terminal
note), and \emph{full-suite with denominator note}. In the 19 September
ledger, three scopes are strictly complete: LAB-Bench public
split~\cite{laurent2024labbench} (E8 replacement), AgentDojo benign
utility~\cite{debenedetti2024agentdojo} (E10 replacement) and
VerilogEval~\cite{liu2023verilogeval} (E11 original). Two further scopes
are qualified, not strictly complete: the
Omni-MATH~\cite{gao2024omnimath} replacement for E14 is terminal-complete
with two judge-excluded rows (the judged denominator 4,426 is canonical;
2,271/4,426 $=0.5131$ judged and 2,271/4,428 $=0.5129$ all-rows are both
shown in Table~\ref{tab:campaign}, and where the source ledger attaches
0.5131 to the 4,428 denominator the judged reading here governs), and
SWE-bench Pro~\cite{deng2025swebenchpro} (E1 original) is full-suite with a
denominator note: its 731-task denominator includes 18 tasks without a
native evaluation. Six later replacement executions (27
September) used a protocol override and do not close their original
contracts; seven original contracts remain externally closed or blocked.
``Fourteen-lane'' therefore names fourteen contracts, not fourteen
completed benchmarks, and the campaign supplies no common cross-lane
estimand. Lane pass rates characterise the null actor of Section~III-E
under each suite's harness (e.g.\ 150/312 VerilogEval extraction failures,
Section~V); they are ledger entries, not capability or training-effect
estimates.

\begin{table}[t]
\caption{Graded results for the trained actor. Each row uses its suite's
native evaluator, metric and denominator (acc.\ = accuracy; ``all rows''
counts the two judge-excluded rows as incorrect); rows are not comparable
and are never aggregated.}
\label{tab:campaign}
\centering
\footnotesize
\setlength{\tabcolsep}{3pt}
\begin{tabular}{@{}llll@{}}
\toprule
Lane, suite (scope) & Metric & Passed/total & Score \\
\midrule
E8 LAB-Bench (repl.) & accuracy & 450/1,967 & 0.2288 \\
E10 AgentDojo (repl.) & benign utility & 88/97 & 0.9072 \\
E11 VerilogEval (orig.) & pass@1 & 129/312 & 0.4135 \\
E14 Omni-MATH (repl.) & acc., judged rows & 2,271/4,426 & 0.5131 \\
E14 Omni-MATH (repl.) & acc., all rows & 2,271/4,428 & 0.5129 \\
E1 SWE-bench Pro (orig.) & pass@1 & 2/731 & 0.0027 \\
\bottomrule
\end{tabular}
\end{table}

\begin{table*}[t]
\caption{E1--E14 lanes: original contract, declared replacement scope, and
terminal state of each scope (19 September ledger; 27 September replacement
executions marked ``exploratory''). Complete = full scope-specific
dispositions. Closed ext.\ = external access closed with a written record.}
\label{tab:lanes}
\centering
\footnotesize
\setlength{\tabcolsep}{3pt}
\begin{tabular}{@{}lllll@{}}
\toprule
Lane & Original contract & State (original) & Replacement scope & State (replacement) \\
\midrule
E1 & SWE-bench Pro~\cite{deng2025swebenchpro} & complete (2/731) & SWE-bench Multilingual~\cite{khandpurswemultilingual} & exploratory \\
E2 & FrontierSWE~\cite{proximal2026frontierswe} & partial (1/17 tasks) & CORE-Bench~\cite{siegel2024corebench} & exploratory \\
E3 & SDAB~\cite{emulated2026sdab} & closed ext. & --- & --- \\
E4 & BankerToolBench~\cite{lau2026bankertoolbench} & closed, partial (1/100 tasks) & --- & --- \\
E5 & APEX-Agents~\cite{vidgen2026apexagents} & partial (7/480 scored) & Tau3 & exploratory \\
E6 & WebBench~\cite{halluminatewebbench} & no result & WebArena~\cite{zhou2024webarena} & exploratory (108 ungraded) \\
E7 & BinaryAudit~\cite{quesma2026binaryaudit} & closed ext. & --- & --- \\
E8 & LifeSciBench~\cite{liu2026lifescibench} & closed ext. & LAB-Bench public~\cite{laurent2024labbench} & complete \\
E9 & MLE-bench~\cite{chan2025mlebench} & partial (40/75 graded) & ML-Dev-Bench~\cite{padigela2025mldevbench} & exploratory \\
E10 & AgentHarm~\cite{andriushchenko2025agentharm} & closed ext. & AgentDojo benign~\cite{debenedetti2024agentdojo} & complete \\
E11 & VerilogEval~\cite{liu2023verilogeval} & complete & --- & --- \\
E12 & AppBench~\cite{afterquery2025appbench} & closed ext. & --- & --- \\
E13 & OpenReward games~\cite{openrewardplatform} & blocked, no record & BALROG~\cite{paglieri2025balrog} & exploratory \\
E14 & FrontierMath~\cite{glazer2024frontiermath} & closed ext. & Omni-MATH~\cite{gao2024omnimath} & complete, terminal note \\
\bottomrule
\end{tabular}
\end{table*}

\section{Discussion and Threats to Validity}

The primary training-effect statement of this paper is the paired rerun:
no trained-versus-base difference was found, so every E1--E14 score must be
read as a ledger entry for a null actor, not as a training outcome. In
detail, the campaign results describe one actor at one revision; no full-suite
base-model comparator exists, so they are not evidence of what the GRPO step
contributed. A separate small paired rerun (26 September) scored the trained
actor and the base model on identical items under one vLLM configuration,
5--151 items per lane; for externally blocked lanes it used public subsets
or declared public substitutes, so it is not a rerun of the original
contracts. No lane showed a trained-versus-base difference (smallest
McNemar $p=0.25$). \emph{Length-bias threat (condensed).} Length bias is
not established and is not tested as a result: a post-hoc roster flag
(reward-based, post-hoc thresholds) is not a length test, and the only
controlled length comparison (Qwen2.5-1.5B-Instruct, GSM8K, three seeds)
was right-censored at a 200-token cap with starting lengths near the cap
(mean length fell, 195.0$\to$180.3 GRPO, 193.9$\to$186.8 Dr.\ GRPO; gain
difference $+0.012$, $t_2=1.61$, $p=0.25$, inconclusive), so this
configuration cannot test unconstrained length inflation. The training run itself was 40 steps with reward zero on
22 of them. Tinker results measure the platform's GRPO implementation,
whose internals are closed. Several scores carry harness limitations (e.g.,
150 of 312 VerilogEval responses failed answer extraction in a run with a
1,024-token cap and thinking enabled; per-sample causes were not recorded),
and the VerilogEval managed sampler path no longer exists,
so a re-run would change the sampler as well as the run. The ledger therefore
supports scope accounting and failure-mode auditing, not reproduction of
every score. These limits are stated with the numbers because hiding them
would overstate the evidence.

\section{Conclusion}
The contribution is provenance and ledger discipline, not attribution: the
uncrossed corpus cannot allocate effects to stack versus algorithm, and no
independent stack attribution or algorithm equivalence is established.
Matched same-stack contrasts were inconclusive at their detectable effect
sizes, ZVF predictive validity is unestablished (ZVF is retained as
bookkeeping over logged rewards), and the fourteen-lane ledger is a
governance case study on a null actor. No controller benefit or
post-training capability gain is established. Evidence-bound practice rules
within these limits: (1) report seed counts with MDEs---at the dispersions
observed here, eight paired seeds still leave MDEs of 0.010--0.015, so
sub-point variant differences at $\le$8 seeds should be reported as
inconclusive, never as equivalence; (2) log per-group reward tensors (or at
minimum ZVF with mean reward, estimator form and $G$) before launch, so
bookkeeping can be recomputed rather than trusted; (3) treat larger $G$ as
buying variance-bearing groups at proportional rollout cost (ZVF
0.838$\to$0.631 for $G=2\to16$ here), not as a free accuracy improvement;
(4) freeze a lane ledger with original/replacement scopes and terminal
states before scoring, so closed and null lanes remain auditable instead of
disappearing.

\appendices
\section{Artifact, Availability and Overlap Statements}
\emph{Artifacts and recompute entry points (as-is).} Repository commit:
N/A---no archival commit, DOI or release tag is bound to this draft, and
none is claimed (AUTHOR\_TODO: bind an archival commit/release before
submission). The following paths exist in the project repository at the
time of writing and reproduce the recomputed numbers: Phase-1 tensors
\texttt{platform\_hybrid/experiments/results/n2\_reward\_tensor\_resume/}
\texttt{\{grpo,aero,gift,areal\}\_s0\_tensors.jsonl} (SHA-256 prefixes
\texttt{475c717f}, \texttt{d3885dd6}, \texttt{6fa60ff8}, \texttt{5c1e6a89})
via \texttt{tools/review\_ledger/PES\_IEEE\_Attribution\_recompute.py}
(output \texttt{review\_2026-10-18/PES\_IEEE\_Attribution/recompute.json});
Fig.~\ref{fig:gu} from \texttt{groupsize\_zvf\_sweep.json} via the same
script; PPO/GRPO per-seed files under
\texttt{modal\_registry/modal\_artifacts/tinkerrl-results/samestack/};
variant contrasts from \texttt{e1\_confirmatory\_wandb\_export.tsv}
(W\&B project \texttt{tinker-rl-lab}, runs \texttt{e1-confirmatory-*},
exported 2026-10-18) via
\texttt{PES\_IEEE\_Attribution\_p11\_recompute.py}; W\&B binning via
\texttt{PES\_IEEE\_Attribution\_zvf\_audit.py}
(\texttt{zvf\_audit\_recompute.json}); ledger code check
\texttt{outputs/verification/LEDGER\_CODE\_CHECK\_2026-09-19.json}
(11/11 PASS). \emph{Non-reproducible by design or by loss:} the 80 drawn
training prompts, consumed-row inventory and executed actor training code
(unrecovered); per-lane receipts behind LAB-Bench and AgentDojo counts
(not bound in this paper); Tinker run IDs, adapter checkpoint hash and the
defunct VerilogEval managed sampler path; closed managed-platform
internals generally. The manuscript alone therefore does not support
independent recomputation of Table~\ref{tab:campaign}.
\emph{Data/code availability (truthful, including absences).} No public
release is claimed. Code and result files exist only in the authors'
project repository as listed above; benchmark lane data remain under their
owners' terms; seven lanes are externally closed. Compute was short-horizon
academic post-training on a managed closed platform (Tinker) and local/TRL
stacks; no costed compute statement or funding/acknowledgement can be
verified from recovered records (AUTHOR\_TODO: supply funding,
acknowledgements, platform-version pins and any conflict-of-interest note
for the closed platform). No venue, organiser or AI-use declaration is
asserted in this draft (AUTHOR\_TODO).
\emph{Reference verification status.} Six 2026-dated grey/vendor sources
used only as lane provenance (\cite{lau2026bankertoolbench},
\cite{vidgen2026apexagents}, \cite{liu2026lifescibench},
\cite{emulated2026sdab}, \cite{quesma2026binaryaudit},
\cite{proximal2026frontierswe}) are unverified as archival publications
and are cited only to name closed/partial lane contracts; they support no
graded result. AUTHOR\_TODO: verify each against its canonical record or
move lane provenance to a URL list, and confirm key/year details for
\cite{thinkingmachines2024tinker}, \cite{rafailov2024direct},
\cite{le2025rlzvp} and the undated \cite{khandpurswemultilingual}.
\emph{Overlap statement.} Sibling manuscripts from the same project share
the underlying archive: the ICICC submission \emph{Receipt-Bound Evaluation
Governance: A Fourteen-Lane Language-Model Case Study} reports the same
E1--E14 ledger as its primary subject; the M.Tech thesis contains the
project-wide record from which these numbers are drawn; two further
siblings are disclosed by title only here---\emph{The Stack Is Part of the
Result: Negative Results on Attributing GRPO-Family Differences} (ACL ARR)
and \emph{Signal Starvation in GRPO: Measuring and Acting on Zero-Variance
Groups} (RAIT). What is unique here is the combination at IEEE-conference
length: tensor-backed ZVF bookkeeping, the same-stack contrast statistics
of Table~\ref{tab:variants}, and the ledger summarised under the single
taxonomy of Section~IV-F. Editors should request the siblings for
comparison; shared tables must not be read as independent confirmations
(AUTHOR\_TODO: finalise per-sibling shared-data/text enumeration with
collaborator authorisation before submission).

\bibliographystyle{IEEEtran}
\bibliography{refs}

\end{document}
